% =====================================================================
%  29.1  Bridges and the bridge tree.
%  Exports: sec:c31-bridges  def:c31-bridge  thm:c31-bridgetest
%  thm:c31-robbins  def:c31-twiedge  lem:c31-lowchar  lst:c31-bridge
%  lst:c31-contract  prob:c31-roads  fig:c31-roads  lst:c31-roads
%  Solved CF 118E "Bertown roads" (2000, L4), live-verified against the CF
%  problem page + problemset API (name "Bertown roads", rating 2000, tags
%  dfs and similar / digraphs / graphs); two official samples: the bridgeless
%  6x8 graph -> a valid strong orientation, the 6x7 graph with bridge 1-4 -> 0.
%  This section OWNS the low-link bridge test and the 2-edge-connected
%  contraction; 29.2-29.5 and the Summit reference them by label.
% =====================================================================
\section{Bridges and the Bridge Tree}\label{sec:c31-bridges}

\begin{previouslybox}[Ch.28 assumed the input was a tree]
Chapter~28 (\chref{c22}) could only compare two trees once you \emph{knew} the
input was acyclic and connected.  But so many contest inputs arrive as a general
graph and the real question is \emph{edge survival}: which edges, if removed,
break the network apart?  Answer that and you learn exactly where the cycles are,
and the acyclic skeleton that remains -- the bridge tree -- is a tree, so
Ch.~12 (\chref{c09}) and Ch.~13 (\chref{c10}) become legal on it.  This section
builds that skeleton in one DFS and spends it on CF 118E, where the question is
dressed up as orienting roads.
\end{previouslybox}

\begin{topicintro}
A \term{bridge} is an edge whose deletion increases the number of connected
components (\(\defref{c31-bridge}\)).  Tarjan's one-pass DFS labels each vertex
with an entry time \(\mathrm{tin}\) and a \emph{low} value -- the smallest
\(\mathrm{tin}\) its subtree can reach using at most one back edge -- and an edge
\(u\!-\!v\) (with \(v\) the child) is a bridge exactly when
\(\mathrm{low}[v]>\mathrm{tin}[u]\) (\(\lemref{c31-lowchar}\)): the child's whole
subtree has no way up except through \(u\!-\!v\).  Delete every bridge and the
remaining pieces are the \term{2-edge-connected components}
(\(\defref{c31-twiedge}\)); add the bridges back as edges between pieces and you
get the \term{bridge tree}.  CF 118E pushes this to its logical end: Robbins'
theorem (\(\thmref{c31-robbins}\)) says a connected graph can be oriented into a
strongly connected digraph \emph{iff} it has no bridge, so the problem collapses
to a single bridge scan plus the DFS edge orientation.

\textbf{Prerequisites.} \(\tin\) from Ch.~5 (\chref{c24}); the connected-component
union of Ch.~17 (\chref{c28}) for the contraction.  No prior exposure to
low-link is assumed -- \(\lemref{c31-lowchar}\) is proved from scratch here.

\textbf{After this section you can}
\begin{itemize}[nosep]
  \item run the low-link DFS and list every bridge in \(\bigO{n+m}\); \hfill\lv{3}
  \item contract the 2-edge-connected components and hand back the bridge tree; \hfill\lv{4}
  \item state and use Robbins' theorem to turn ``orient for strong connectivity'' into ``is there a bridge?''; \hfill\lv{4}
  \item solve CF 118E -- orient or print \(0\), with the DFS-tree orientation. \hfill\lv{4}
\end{itemize}
\end{topicintro}

% ---------------------------------------------------------------
\subsection{The low-link bridge test}
% ---------------------------------------------------------------

\begin{definition}{Bridge and 2-edge-connected component}{c31-bridge}
An edge \(e\) of an undirected graph \(G\) is a \term{bridge} if
\(G\setminus\{e\}\) has strictly more connected components than \(G\).  Two
vertices are \(2\)\term{-edge-equivalent} if no bridge separates them (equivalently
there are two edge-disjoint paths between them); the equivalence classes are the
\term{2-edge-connected components} (\defref{c31-twiedge} restates this as the
bridge-free components).
\end{definition}

\begin{definition}{2-edge-connected components}{c31-twiedge}
Remove every bridge of \(G\); each maximal connected remaining subgraph is a
2-edge-connected component.  Contract each to a single node and keep the bridges
as edges between nodes: the result is a forest, and it is a tree when \(G\) is
connected -- the \term{bridge tree}.  Its nodes are exactly the classes of
\(\defref{c31-bridge}\).
\end{definition}

\begin{lemma}{Low-link characterisation of a bridge}{c31-lowchar}
Run a DFS from an arbitrary root, timestamp each vertex \(v\) with
\(\mathrm{tin}[v]\), and set
\[
  \mathrm{low}[v] \;=\; \min\Bigl\{\, \mathrm{tin}[v],\;\; \min_{(v,x)\text{ back}}\mathrm{tin}[x],\;\; \min_{v\to c\text{ tree}}\mathrm{low}[c] \Bigr\}.
\]
For a tree edge \(u\!-\!v\) (rooted so \(v\) is the child), \(u\!-\!v\) is a
bridge \emph{iff} \(\mathrm{low}[v]>\mathrm{tin}[u]\).
\end{lemma}

\begin{proofbox}
\((\Rightarrow)\) Suppose \(u\!-\!v\) is a bridge but \(\mathrm{low}[v]\le\mathrm{tin}[u]\).
By the definition of low there is a path from inside \(v\)'s subtree to some
ancestor \(a\) of \(v\) (possibly \(u\)) that uses the single back edge
\((\cdot,a)\) and no tree edge above \(u\); that path together with the tree
path \(u\leadsto v\) avoids the edge \(u\!-\!v\), giving a second route between
\(u\) and \(v\), so deleting \(u\!-\!v\) keeps them connected -- contradiction.
\((\Leftarrow)\) If \(\mathrm{low}[v]>\mathrm{tin}[u]\), every edge with one end
in \(v\)'s subtree and the other outside is a tree edge, and the only such edge
is \(u\!-\!v\) itself (a back edge reaching \(\le\mathrm{tin}[u]\) would force
\(\mathrm{low}[v]\le\mathrm{tin}[u]\)).  Hence no alternate route leaves the
subtree, and cutting \(u\!-\!v\) separates it from \(u\). \(\Box\)
\end{proofbox}

\begin{theorem}{Robbins' theorem}{c31-robbins}
A connected undirected graph admits an orientation making the digraph strongly
connected \emph{if and only if} it has no bridge.
\end{theorem}

\begin{proofbox}
\emph{(Only if.)} If \(u\!-\!v\) is a bridge it is the unique connection between
the two components \(A\ni u\), \(B\ni v\); orient it \(u\to v\) and no vertex of
\(B\) can return to \(A\).  \emph{(If.)} If there is no bridge, every edge lies on
a simple cycle.  Root a DFS at any \(r\).  Orient each tree edge \(p\to c\)
downward and each back edge from the deeper endpoint up to the shallower one.
Then \(r\) reaches every vertex along tree edges, and every vertex reaches \(r\):
its subtree root has a back edge upward (there are no bridges, so
\(\mathrm{low}[\text{child}]\le\mathrm{tin}[\text{parent}]\) and some back edge
escapes).  Composing \(x\to r\to y\) gives strong connectivity. \(\Box\)
\end{proofbox}

\begin{lstlisting}[caption={Low-link bridges plus 2-edge-connected contraction. One DFS fills tin and low, marks a bridge when low[child] is above tin[parent], then a second pass labels components by skipping bridges. ASCII only, identifiers without underscores}, label={lst:c31-bridge}]
#include <bits/stdc++.h>
using namespace std;
const int MAXN=100005;
int n,m,timer=0; vector<pair<int,int> > adj[MAXN];
int tin[MAXN], low[MAXN], comp[MAXN], isBridge[MAXN];
void dfsBridge(int u,int pe){
  tin[u]=low[u]=++timer;
  for(auto &pr : adj[u]){
    int v=pr.second, id=pr.first;
    if(id==pe) continue;                    // skip the SAME edge, not parent vertex
    if(tin[v]){ low[u]=min(low[u], tin[v]); }   // back edge to ancestor/visited
    else { dfsBridge(v,id); low[u]=min(low[u],low[v]);
           if(low[v]>tin[u]) isBridge[id]=1; }  // the bridge test
  }
}
void dfsComp(int u,int c){                   // contract: never cross a bridge
  comp[u]=c;
  for(auto &pr : adj[u]) if(!comp[pr.second] && !isBridge[pr.first]) dfsComp(pr.second,c);
}
// usage: for all u if(!tin[u]) dfsBridge(u,-1); then c=0 for all u if(!comp[u]) dfsComp(u,++c);
\end{lstlisting}

\begin{insightbox}[Skip the edge id, not the parent vertex]
The single bug that ruins most low-link code is writing \code{if(v==parent) continue;}.
With multi-edges that is wrong: two parallel edges \(u\!-\!v\) form a 2-cycle, so
neither is a bridge, yet skipping \emph{the vertex} \(u\) discards both.  Track the
\emph{edge id} (\(\coderef{c31-bridge}\)'s \code{pe}) and skip only the exact edge
you came down; the parallel edge then correctly lowers \(\mathrm{low}\).
\end{insightbox}

\begin{complexitybox}[Bridges and contraction]
One DFS computes \(\mathrm{tin}/\mathrm{low}\) and every bridge in
\(\bigO{n+m}\); the contraction DFS is another \(\bigO{n+m}\).  Space
\(\bigO{n+m}\).  The bridge tree has at most \(n\) nodes and \(n-1\) bridge
edges, so anything Ch.~12/13 do on a tree costs \(\bigO{n}\) more on top.
\end{complexitybox}

% ---------------------------------------------------------------
\subsection{Problem Lab: CF 118E Bertown roads -- orient or give up}
% ---------------------------------------------------------------

\solvedproblem{c31-roads}{Bertown roads}{CF 118E}
  {https://codeforces.com/problemset/problem/118/E}{CF 2000 L4}
  {Robbins' theorem (\(\thmref{c31-robbins}\)): orientable iff bridgeless; the DFS that
   finds bridges (\(\lemref{c31-lowchar}\), \(\coderef{c31-bridge}\)) simultaneously produces
   the strong orientation -- tree edges downward, back edges up toward the ancestor}

\begin{problemstatement}[Bertown roads (statement, constraints and both samples verified against the live CF page)]
Bertown is a network of \(n\) junctions (\(2\le n\le 1000\)) joined by \(m\) two-way roads
(\(1\le m\le \min\!\bigl(n(n-1)/2,\,2000\bigr)\)); the graph is connected and has no
multi-edges or self-loops.  The mayor wants every road made \textbf{one-way} so that from
\emph{any} junction you can drive to \emph{any other} junction -- the resulting directed
graph must be strongly connected.

\textbf{Task.} If no such one-way assignment exists, print a single line \texttt{0}.
Otherwise print \(m\) lines, each ``\(x\ y\)'': the road \(\{x,y\}\) from the input becomes
directed \(x\to y\).  Every input road must appear exactly once in some orientation.  Any
valid answer is accepted.  Time limit 2 seconds.

\textbf{Samples (both official).} \#1: \(n{=}6,m{=}8\) with roads
\(1\!-\!2,2\!-\!3,3\!-\!1,4\!-\!5,5\!-\!6,6\!-\!4,1\!-\!4,2\!-\!5\) \(\to\) a legal
orientation such as \texttt{1 2 / 2 3 / 3 1 / 5 4 / 6 5 / 4 6 / 4 1 / 5 2}.  \#2:
\(n{=}6,m{=}7\) with roads \(1\!-\!2,2\!-\!3,3\!-\!1,4\!-\!5,5\!-\!6,6\!-\!4,1\!-\!4\)
\(\to\) \texttt{0}.
\end{problemstatement}

\begin{figure}[H]
\centering
\begin{tikzpicture}[diagram, scale=0.95, transform shape,
   tnode/.style={circle, draw=primaryblue!70, fill=primaryblue!8, line width=0.7pt,
                 minimum size=7mm, font=\scriptsize\bfseries, inner sep=1pt},
   every edge/.style={draw, -Latex, line width=0.9pt}]
  \node[tnode] (v1) at (0,1.6) {1};   \node[tnode] (v2) at (-1.4,0) {2};   \node[tnode] (v3) at (0.6,0) {3};
  \node[tnode] (v4) at (3.2,1.6) {4}; \node[tnode] (v5) at (3.2,0) {5};    \node[tnode] (v6) at (4.6,0.8) {6};
  \draw (v1)--(v2) (v2)--(v3) (v3)--(v1) (v4)--(v5) (v5)--(v6) (v6)--(v4) (v1)--(v4) (v2)--(v5);
  \node[annot, font=\scriptsize] at (1.7,-0.9) {sample 1: two triangles joined by two rungs -> no bridge};
\end{tikzpicture}
\caption{Sample 1 of CF 118E.  The two rungs \(1\!-\!4\) and \(2\!-\!5\) make every edge
lie on a cycle, so the graph is bridgeless and Robbins' theorem guarantees a strong
orientation.  Sample 2 deletes the rung \(2\!-\!5\), leaving \(1\!-\!4\) a bridge between
the two triangles -- the answer flips to \(0\).}
\label{fig:c31-roads}
\end{figure}

\begin{attemptgate}[Attempt it yourself first -- \lv{4}, give yourself 25 minutes]
Before coding: what property of the undirected graph makes the task
\emph{impossible}, and if the graph has that property absent, which orientation of a DFS
tree is forced?
\peek{If some road is a bridge, every journey between the two sides must use it, so once
you orient it one way you can never return -- a bridge makes the task impossible.  This is
Robbins' theorem (\(\thmref{c31-robbins}\)): orientable iff bridgeless.}
\nudge{So run the bridge DFS of \(\coderef{c31-bridge}\); if it flags any bridge, print \(0\)
and stop.  Otherwise you already have a DFS tree -- reuse it: orient tree edges
parent\(\to\)child and back edges from the deeper vertex up to the ancestor.}
\fullhint{Root a DFS at vertex \(1\).  While exploring, classify each road: a \emph{tree}
edge is one that discovers a new vertex, a \emph{back} edge is one whose other end already
has a \(\mathrm{tin}\).  Emit tree edges as \(u\to v\) (down) and back edges as
\(u\to v\) where \(\mathrm{tin}[v]<\mathrm{tin}[u]\) (up toward the ancestor).  Because
there is no bridge, every tree edge \(u\!-\!v\) satisfies \(\mathrm{low}[v]\le\mathrm{tin}[u]\),
so a back edge escapes each subtree: from any vertex you can climb to \(1\) (tree edges let
\(1\) reach you, back edges plus tree edges let you return), giving strong connectivity.
Track the parent \emph{edge id}, not the parent vertex, so you do not skip a parallel edge
(none exist here, but make the habit).}
\end{attemptgate}

\solstep{Understanding}
CF 118E is the bridge test wearing a costume.  ``Can I orient every edge so traffic
circulates both ways?'' is precisely ``is there an unavoidable road?'' -- a bridge is exactly
an edge that all one-directional traffic must funnel through.  Robbins' theorem is the biconditional,
and it is \emph{constructive}: the DFS that certifies ``no bridge'' also emits the witness
orientation, tree edges down and back edges up.  So you run one DFS and either print \(0\) or
print the orientation it naturally produced -- there is nothing to solve after the scan.

\solstep{Key Observations}
\begin{observation}
A bridge, once oriented, is a one-way door between its two sides, so a strongly connected
orientation can never exist while a bridge does.  Conversely, in the DFS orientation of
\(\thmref{c31-robbins}\)'s proof, a vertex returns to the root along the \emph{back edge}
that escapes its subtree -- which is guaranteed to exist exactly because
\(\mathrm{low}[v]\le\mathrm{tin}[u]\) says no tree edge is a bridge.
\end{observation}
\begin{observation}
The orientation is not unique -- any witness is accepted.  The one the DFS emits is the
cheapest to produce (no extra pass): the direction of road \(\{u,v\}\) is decided by which
endpoint was visited first (tree edge) or by the \(\mathrm{tin}\) order (back edge), so you
just print in the DFS and never revisit.
\end{observation}

\solstep{The Approach Ladder}
\begin{center}\renewcommand{\arraystretch}{1.25}\setlength{\tabcolsep}{3pt}
\begin{tabular}{@{}c p{0.26\linewidth} p{0.30\linewidth} l c@{}}
\toprule
\textbf{\#} & \textbf{Approach} & \textbf{Idea} & \textbf{Time} & \textbf{Verdict}\\
\midrule
1 & \lv{2}\ try all \(2^m\) orientations & brute force sign per edge & \(\bigO{2^m}\) & \tle \\
2 & \lv{3}\ SCC-check a guessed orientation & pick arbitrary, verify & \(\bigO{m(n+m)}\) & \textbf{WA} \\
3 & \lv{4}\ \textbf{bridge DFS + free orientation} & Robbins: orientable iff bridgeless & \(\bigO{n+m}\) & \ac \\
\bottomrule
\end{tabular}\end{center}
Rung~2 is tempting (orient greedily, run Tarjan SCC, retry) but has no principled stopping
rule and is quadratic in the worst case.  Rung~3 is one DFS: it both decides possibility and
emits the witness, because the DFS-tree orientation \emph{is} strong exactly when no bridge
was flagged.

\solstep{Solution}
Read the graph with edge ids.  Run the bridge DFS of \(\coderef{c31-bridge}\) from vertex
\(1\); the same recursion decides each road's direction: when it first reaches \(v\) from \(u\)
emit \(u\to v\); when it sees an already-visited \(v\) with \(\mathrm{tin}[v]<\mathrm{tin}[u]\)
emit \(u\to v\) (the back edge, upward).  If any edge was flagged a bridge, print \(0\);
otherwise print the stored orientations, one road per line.

\solstep{Implementation}
\begin{lstlisting}[caption={CF 118E. One DFS finds bridges AND emits the strong orientation: tree edges downward, back edges upward toward the ancestor. Print 0 if any bridge. ASCII only, identifiers without underscores}, label={lst:c31-roads}]
#include <bits/stdc++.h>
using namespace std;
const int MAXN=1005;
int n,m,timer=0,hasBridge=0;
vector<pair<int,int> > adj[MAXN];
int tin[MAXN], low[MAXN], dirA[MAXN], dirB[MAXN], seenEdge[MAXN];
void dfs(int u,int pe){
  tin[u]=low[u]=++timer;
  for(auto &pr : adj[u]){
    int v=pr.second, id=pr.first;
    if(id==pe) continue;
    if(tin[v]){
      low[u]=min(low[u], tin[v]);
      if(tin[v]<tin[u] && !seenEdge[id]){ seenEdge[id]=1; dirA[id]=u; dirB[id]=v; } // back: u->v up
    } else {
      dfs(v,id); low[u]=min(low[u], low[v]);
      if(low[v]>tin[u]) hasBridge=1;             // a bridge -> impossible
      seenEdge[id]=1; dirA[id]=u; dirB[id]=v;    // tree: u->v down
    }
  }
}
int main(){
  scanf("%d%d",&n,&m);
  for(int i=1;i<=m;i++){ int u,v; scanf("%d%d",&u,&v);
    adj[u].push_back({i,v}); adj[v].push_back({i,u}); }
  dfs(1,-1);
  if(hasBridge){ printf("0\n"); return 0; }
  for(int i=1;i<=m;i++) printf("%d %d\n",dirA[i],dirB[i]);
}
\end{lstlisting}

\solstep{Trace}
\begin{tracebox}[Sample 2 in full: the bridge that decides everything]
Run \(\coderef{c31-roads}\) from vertex \(1\).  DFS reaches \(2,3\) and closes the back edge
\(3\!-\!1\) (\(\mathrm{low}[3]{=}1\le \mathrm{tin}[2]\)), so the triangle \(B_1\) has no
bridge; from \(1\) it crosses \(1\!-\!4\) and covers \(5,6\) with back edge \(6\!-\!4\).  When
the recursion returns to \(1\) along the tree edge \(1\!-\!4\), the child \(4\) has
\(\mathrm{low}[4]=\mathrm{tin}[4]=4\) (the second triangle cannot reach \(\mathrm{tin}[1]=1\)
because \(1\!-\!4\) is its only exit).  Since \(\mathrm{low}[4]=4>\mathrm{tin}[1]=1\), edge
\(1\!-\!4\) is a bridge: \(\mathit{hasBridge}=1\).  Print \(\boxed{0}\) -- matches the official
output.  \(\Box\)
\end{tracebox}
\begin{tracebox}[Sample 1: no bridge, so the DFS orientation is the answer]
Run \(\coderef{c31-roads}\) from vertex \(1\) (\(\mathrm{tin}[1]{=}1\)).  Take tree edges
\(1\!-\!2,2\!-\!3,2\!-\!5,5\!-\!4,4\!-\!6\), so the DFS tree already spans all six vertices
and the road \(1\!-\!4\) is a \emph{back} edge (vertex \(4\) was reached through \(5\), not
through \(1\)).  The back edges are \(3\!-\!1\), \(4\!-\!1\) and \(6\!-\!5\), each dropping a
subtree's \(\mathrm{low}\) to \(1\).  Now test every tree edge: under \(2\!-\!3\) the child
\(3\) escapes via \(3\!-\!1\) so \(\mathrm{low}[3]{=}1\le\mathrm{tin}[2]\); under
\(2\!-\!5\) the whole \(\{5,4,6\}\) subtree escapes via \(4\!-\!1\) so \(\mathrm{low}[5]{=}1
\le\mathrm{tin}[2]\); likewise \(\mathrm{low}[4],\mathrm{low}[6]\le\mathrm{tin}\) of their
parents.  No tree edge satisfies \(\mathrm{low}[\text{child}]>\mathrm{tin}[\text{parent}]\),
so \(\mathit{hasBridge}=0\).  The directions the DFS stored -- tree edges downward
(\(1\!\to\!2,2\!\to\!3,2\!\to\!5,5\!\to\!4,4\!\to\!6\)) and back edges upward
(\(3\!\to\!1,4\!\to\!1,6\!\to\!5\)) -- form one strongly connected digraph (from any vertex
climb a back edge toward \(1\), then descend tree edges to any target), and the judge accepts
it: the stored lines \(\boxed{\texttt{1 2 / 2 3 / 2 5 / 5 4 / 4 6 / 3 1 / 4 1 / 6 5}}\) (the
exact set depends on the DFS order). \(\Box\)
\end{tracebox}

\begin{warningbox}[Three ways this goes wrong]
\begin{enumerate}[nosep,leftmargin=1.6em]
  \item \textbf{Skipping the parent vertex instead of the parent edge.} Harmless on
        CF 118E (no multi-edges) but fatal in general and a habit that bites on
        \(\secref{c31-block}\); pass the edge id.
  \item \textbf{Emitting an edge twice.} A back edge \((u,v)\) is seen from both ends;
        orient and print it only when \(\mathrm{tin}[v]<\mathrm{tin}[u]\) (the
        \code{seenEdge} guard in \(\coderef{c31-roads}\)), or you print a road twice.
  \item \textbf{Re-rooting resets \(\mathrm{low}\).} A bridge is DFS-root independent;
        running a fresh DFS per component is fine, but re-timestamping inside one search
        makes every \(\mathrm{low}\) wrong -- start each vertex exactly once.
\end{enumerate}
\end{warningbox}

\begin{hardeningbox}
\begin{itemize}[nosep,leftmargin=1.4em]
  \item \textbf{Recursion depth.} \(n\le1000\) here is safe, but on the contraction
        (\(\coderef{c31-bridge}\), \(n\le10^5\)) an iterated DFS or an increased stack is
        required -- the classic segfault-on-the-chain test.
  \item \textbf{Disconnected input.} If the graph is not connected the problem is moot, but
        your bridge loop must still start a DFS in every unvisited component; CF 118E
        promises connectivity, generic bridge code should not.
  \item \textbf{Verify a YES.} The judge accepts any strong orientation; when you want to
        self-check, run Tarjan SCC on your output and confirm a single component -- cheap
        insurance while developing.
\end{itemize}
\end{hardeningbox}

\begin{recallbox}[29.1 in one breath]
\emph{A bridge is an unavoidable edge; \(\mathrm{low}[v]>\mathrm{tin}[u]\) finds every one in
\(\bigO{n+m}\) (\(\lemref{c31-lowchar}\)).  Delete bridges and contract, and the
2-edge-connected components glue into the \textbf{bridge tree} (\(\defref{c31-twiedge}\)).
Robbins' theorem (\(\thmref{c31-robbins}\)) makes CF 118E a one-DFS problem: orientable iff
bridgeless, and the DFS tree itself is the strong orientation (tree edges down, back edges up)
-- print \(0\) the instant a bridge appears (\(\probref{c31-roads}\)).  Skip the parent
\emph{edge}, and \(\mathcal{O}(n+m)\) you hold a tree to run Ch.12/13 on.}
\selftest{
  \item Why does a bridge make CF 118E impossible, in one sentence? \emph{(Once oriented it is a one-way door between its two sides, so no path returns -- Robbins' theorem \(\thmref{c31-robbins}\).)}
  \item State the bridge test for a tree edge \(u\!-\!v\). \emph{(\(\mathrm{low}[v]>\mathrm{tin}[u]\): the child's subtree cannot reach \(u\) or above without the edge \(\lemref{c31-lowchar}\).)}
  \item What are the components you contract, and what do the bridges become? \emph{(The 2-edge-connected components become nodes; the bridges become the edges of the bridge tree \(\defref{c31-twiedge}\).)}
}
\end{recallbox}
